% =====================================================================
%  EJEMPLO — Ficha Volere del RF-01 (Registro de solicitud de tutoria)
%  Ubicacion: docs/requisitos/especificacion/EJEMPLO-RF01-VOLERE.tex
%  Uso: reemplazar los contenidos con los propios requisitos del equipo.
% =====================================================================

\documentclass[11pt,a4paper]{article}
% =====================================================================
%  PREAMBULO LATEX COMPARTIDO - GUIAS DEL PFC INGENIERIA DE REQUISITOS
%  Ubicacion: docs/entregas/preamble_ingreq.tex
%  Uso: cada guia hace % =====================================================================
%  PREAMBULO LATEX COMPARTIDO - GUIAS DEL PFC INGENIERIA DE REQUISITOS
%  Ubicacion: docs/entregas/preamble_ingreq.tex
%  Uso: cada guia hace % =====================================================================
%  PREAMBULO LATEX COMPARTIDO - GUIAS DEL PFC INGENIERIA DE REQUISITOS
%  Ubicacion: docs/entregas/preamble_ingreq.tex
%  Uso: cada guia hace \input{preamble_ingreq} en su cabecera.
% =====================================================================

\usepackage[utf8]{inputenc}
\usepackage[T1]{fontenc}
\usepackage[spanish,es-tabla,es-noshorthands]{babel}
\usepackage{lmodern}
\usepackage{microtype}
\usepackage[a4paper,margin=2.5cm,headheight=14.5pt]{geometry}
\usepackage{setspace}\onehalfspacing
\usepackage{parskip}

\usepackage[table]{xcolor}
\definecolor{uteqverde}{RGB}{0,102,51}
\definecolor{uteqdorado}{RGB}{204,153,0}
\definecolor{grisfila}{RGB}{245,245,245}
\definecolor{goldclaro}{RGB}{252,245,220}
\definecolor{verdeclaro}{RGB}{235,247,238}
\definecolor{textogris}{RGB}{60,60,60}
\definecolor{nivelaus}{RGB}{176,42,42}
\definecolor{codebg}{RGB}{248,248,248}

\usepackage{amsmath,amssymb}
\usepackage{graphicx}
\usepackage{tikz}
\usetikzlibrary{arrows.meta,positioning,shapes.geometric,fit,calc}
\usepackage{fancyhdr}
\usepackage[most]{tcolorbox}

\newtcolorbox{cajadefinicion}[1]{
    colback=uteqverde!5,
    colframe=uteqverde,
    fonttitle=\bfseries,
    coltitle=white,
    title={#1},
    boxrule=0.6pt,
    arc=3pt,
    breakable
}

\newtcolorbox{cajaentregable}[1]{
    colback=uteqdorado!8,
    colframe=uteqdorado!80!black,
    fonttitle=\bfseries,
    coltitle=white,
    title={#1},
    boxrule=0.6pt,
    arc=3pt,
    breakable
}

\newtcolorbox{cajacritica}[1]{
    colback=red!5,
    colframe=red!70!black,
    fonttitle=\bfseries,
    coltitle=white,
    title={#1},
    boxrule=0.6pt,
    arc=3pt,
    breakable
}

\newtcolorbox{cajarepo}[1]{
    colback=blue!4,
    colframe=blue!60!black,
    fonttitle=\bfseries,
    coltitle=white,
    title={#1},
    boxrule=0.6pt,
    arc=3pt,
    breakable
}

\newtcolorbox{cajaconcepto}[1]{
    colback=verdeclaro,
    colframe=uteqverde!70!black,
    fonttitle=\bfseries,
    coltitle=white,
    title={#1},
    boxrule=0.5pt,
    arc=2pt,
    breakable
}

\usepackage{listings}
\lstset{
    basicstyle=\ttfamily\footnotesize,
    backgroundcolor=\color{codebg},
    frame=single,
    breaklines=true,
    columns=fullflexible,
    keepspaces=true
}

\usepackage{array}
\usepackage{booktabs}
\usepackage{longtable}
\usepackage{tabularx}
\usepackage{multirow}

\usepackage[hidelinks,bookmarks=true,bookmarksnumbered=true]{hyperref}
\usepackage{enumitem}\setlist{itemsep=2pt,topsep=2pt}

\newcolumntype{C}[1]{>{\centering\arraybackslash}p{#1}}
\newcolumntype{L}[1]{>{\raggedright\arraybackslash}p{#1}}
 en su cabecera.
% =====================================================================

\usepackage[utf8]{inputenc}
\usepackage[T1]{fontenc}
\usepackage[spanish,es-tabla,es-noshorthands]{babel}
\usepackage{lmodern}
\usepackage{microtype}
\usepackage[a4paper,margin=2.5cm,headheight=14.5pt]{geometry}
\usepackage{setspace}\onehalfspacing
\usepackage{parskip}

\usepackage[table]{xcolor}
\definecolor{uteqverde}{RGB}{0,102,51}
\definecolor{uteqdorado}{RGB}{204,153,0}
\definecolor{grisfila}{RGB}{245,245,245}
\definecolor{goldclaro}{RGB}{252,245,220}
\definecolor{verdeclaro}{RGB}{235,247,238}
\definecolor{textogris}{RGB}{60,60,60}
\definecolor{nivelaus}{RGB}{176,42,42}
\definecolor{codebg}{RGB}{248,248,248}

\usepackage{amsmath,amssymb}
\usepackage{graphicx}
\usepackage{tikz}
\usetikzlibrary{arrows.meta,positioning,shapes.geometric,fit,calc}
\usepackage{fancyhdr}
\usepackage[most]{tcolorbox}

\newtcolorbox{cajadefinicion}[1]{
    colback=uteqverde!5,
    colframe=uteqverde,
    fonttitle=\bfseries,
    coltitle=white,
    title={#1},
    boxrule=0.6pt,
    arc=3pt,
    breakable
}

\newtcolorbox{cajaentregable}[1]{
    colback=uteqdorado!8,
    colframe=uteqdorado!80!black,
    fonttitle=\bfseries,
    coltitle=white,
    title={#1},
    boxrule=0.6pt,
    arc=3pt,
    breakable
}

\newtcolorbox{cajacritica}[1]{
    colback=red!5,
    colframe=red!70!black,
    fonttitle=\bfseries,
    coltitle=white,
    title={#1},
    boxrule=0.6pt,
    arc=3pt,
    breakable
}

\newtcolorbox{cajarepo}[1]{
    colback=blue!4,
    colframe=blue!60!black,
    fonttitle=\bfseries,
    coltitle=white,
    title={#1},
    boxrule=0.6pt,
    arc=3pt,
    breakable
}

\newtcolorbox{cajaconcepto}[1]{
    colback=verdeclaro,
    colframe=uteqverde!70!black,
    fonttitle=\bfseries,
    coltitle=white,
    title={#1},
    boxrule=0.5pt,
    arc=2pt,
    breakable
}

\usepackage{listings}
\lstset{
    basicstyle=\ttfamily\footnotesize,
    backgroundcolor=\color{codebg},
    frame=single,
    breaklines=true,
    columns=fullflexible,
    keepspaces=true
}

\usepackage{array}
\usepackage{booktabs}
\usepackage{longtable}
\usepackage{tabularx}
\usepackage{multirow}

\usepackage[hidelinks,bookmarks=true,bookmarksnumbered=true]{hyperref}
\usepackage{enumitem}\setlist{itemsep=2pt,topsep=2pt}

\newcolumntype{C}[1]{>{\centering\arraybackslash}p{#1}}
\newcolumntype{L}[1]{>{\raggedright\arraybackslash}p{#1}}
 en su cabecera.
% =====================================================================

\usepackage[utf8]{inputenc}
\usepackage[T1]{fontenc}
\usepackage[spanish,es-tabla,es-noshorthands]{babel}
\usepackage{lmodern}
\usepackage{microtype}
\usepackage[a4paper,margin=2.5cm,headheight=14.5pt]{geometry}
\usepackage{setspace}\onehalfspacing
\usepackage{parskip}

\usepackage[table]{xcolor}
\definecolor{uteqverde}{RGB}{0,102,51}
\definecolor{uteqdorado}{RGB}{204,153,0}
\definecolor{grisfila}{RGB}{245,245,245}
\definecolor{goldclaro}{RGB}{252,245,220}
\definecolor{verdeclaro}{RGB}{235,247,238}
\definecolor{textogris}{RGB}{60,60,60}
\definecolor{nivelaus}{RGB}{176,42,42}
\definecolor{codebg}{RGB}{248,248,248}

\usepackage{amsmath,amssymb}
\usepackage{graphicx}
\usepackage{tikz}
\usetikzlibrary{arrows.meta,positioning,shapes.geometric,fit,calc}
\usepackage{fancyhdr}
\usepackage[most]{tcolorbox}

\newtcolorbox{cajadefinicion}[1]{
    colback=uteqverde!5,
    colframe=uteqverde,
    fonttitle=\bfseries,
    coltitle=white,
    title={#1},
    boxrule=0.6pt,
    arc=3pt,
    breakable
}

\newtcolorbox{cajaentregable}[1]{
    colback=uteqdorado!8,
    colframe=uteqdorado!80!black,
    fonttitle=\bfseries,
    coltitle=white,
    title={#1},
    boxrule=0.6pt,
    arc=3pt,
    breakable
}

\newtcolorbox{cajacritica}[1]{
    colback=red!5,
    colframe=red!70!black,
    fonttitle=\bfseries,
    coltitle=white,
    title={#1},
    boxrule=0.6pt,
    arc=3pt,
    breakable
}

\newtcolorbox{cajarepo}[1]{
    colback=blue!4,
    colframe=blue!60!black,
    fonttitle=\bfseries,
    coltitle=white,
    title={#1},
    boxrule=0.6pt,
    arc=3pt,
    breakable
}

\newtcolorbox{cajaconcepto}[1]{
    colback=verdeclaro,
    colframe=uteqverde!70!black,
    fonttitle=\bfseries,
    coltitle=white,
    title={#1},
    boxrule=0.5pt,
    arc=2pt,
    breakable
}

\usepackage{listings}
\lstset{
    basicstyle=\ttfamily\footnotesize,
    backgroundcolor=\color{codebg},
    frame=single,
    breaklines=true,
    columns=fullflexible,
    keepspaces=true
}

\usepackage{array}
\usepackage{booktabs}
\usepackage{longtable}
\usepackage{tabularx}
\usepackage{multirow}

\usepackage[hidelinks,bookmarks=true,bookmarksnumbered=true]{hyperref}
\usepackage{enumitem}\setlist{itemsep=2pt,topsep=2pt}

\newcolumntype{C}[1]{>{\centering\arraybackslash}p{#1}}
\newcolumntype{L}[1]{>{\raggedright\arraybackslash}p{#1}}


\begin{document}

\section*{Ejemplo aplicado --- RF-01 en formato Volere}

\begin{table}[htbp]
\centering
\caption{Ficha Volere del requisito funcional RF-01 (tomado del proyecto modelo, Tabla D1.1).}
\label{tab:rf01}
\rowcolors{2}{grisfila}{white}
\begin{tabularx}{\textwidth}{L{3.7cm}X}
\toprule
\rowcolor{uteqverde!25}
\textbf{Campo} & \textbf{Contenido}\\
\midrule
ID del Requisito             & RF-01\\
Nombre                       & Registro de solicitud de tutor\'ia\\
Tipo                         & Funcional\\
Descripci\'on                & El estudiante deber\'a registrar una solicitud de tutor\'ia
                               seleccionando la asignatura, el tema o motivo, la hora de
                               disponibilidad, su preferencia para el tipo (individual
                               o grupal), la modalidad (presencial o virtual) y cargar
                               un archivo cuando sea necesario.\\
Rationale (Justificaci\'on)  & Es la funcionalidad central para formalizar las solicitudes
                               de tutor\'ia y garantizar trazabilidad.\\
Criterios de aceptaci\'on    & \begin{itemize}[leftmargin=*,nosep]
                                 \item El formulario valida campos obligatorios.
                                 \item La solicitud queda registrada en la base de datos.
                                 \item Se genera notificaci\'on al docente correspondiente.
                               \end{itemize}\\
Prioridad                    & Alta\\
Fuente                       & Encuesta estudiantil (n=40, junio 2025)\\
Restricciones                & Solo se pueden solicitar tutor\'ias de asignaturas
                               activas en el ciclo acad\'emico vigente.\\
Responsable / Actor          & Estudiante\\
Interesados secundarios      & Docente, Coordinaci\'on acad\'emica\\
Casos de uso asociados       & CU-01 (Registrar solicitud de tutor\'ia)\\
Historias de usuario         & HU-06 (Como estudiante, quiero registrar una solicitud
                               de tutor\'ia\ldots)\\
Priorizaci\'on MoSCoW        & \textbf{Must}\\
Priorizaci\'on Kano          & \textbf{Imprescindible} (insatisfactor)\\
M\'etodo de verificaci\'on   & Inspecci\'on del formulario y prueba de aceptaci\'on Gherkin\\
Riesgo                       & Alto (impacto en el flujo principal del sistema)\\
Estado                       & Aprobado\\
\bottomrule
\end{tabularx}
\end{table}

\subsection*{Cumplimiento de INCOSE v4 (autoevaluaci\'on)}

Este requisito cumple con las siguientes caracter\'isticas individuales de INCOSE
Guide to Writing Requirements v4:

\begin{itemize}[leftmargin=*]
    \item \textbf{Necesario}: sin este requisito, el proceso no puede iniciarse.
    \item \textbf{Apropiado}: est\'a al nivel de detalle correcto para su categor\'ia (funcional).
    \item \textbf{No ambiguo}: los t\'erminos ``solicitud'', ``asignatura'', ``tema'', ``modalidad'', ``tipo'' est\'an definidos en el glosario.
    \item \textbf{Completo}: incluye todos los elementos necesarios para su implementaci\'on.
    \item \textbf{Singular}: expresa una \'unica funcionalidad (registrar una solicitud).
    \item \textbf{Factible}: es t\'ecnicamente realizable con la pila propuesta.
    \item \textbf{Verificable}: los criterios de aceptaci\'on son observables.
    \item \textbf{Correcto}: refleja fielmente la necesidad expresada en la encuesta.
    \item \textbf{Conforme}: sigue el patr\'on 29148 (sujeto expl\'icito, verbo modal ``deber\'a'').
    \item \textbf{Expresado positivamente}: no usa ``no''.
    \item \textbf{Apto para la abstracci\'on}: es independiente del framework de implementaci\'on.
    \item \textbf{Consistente}: no entra en conflicto con otros requisitos.
    \item \textbf{Comparable}: puede compararse con requisitos an\'alogos de otros sistemas.
    \item \textbf{Modificable}: puede refinarse en iteraciones futuras sin afectar la estructura.
    \item \textbf{Permitido}: no viola normativa institucional ni legal.
\end{itemize}

\textbf{Total de caracter\'isticas cumplidas}: 15/15 (100\,\%).

\end{document}
